\section{模型的评价}

\subsection{模型优点}

本文以十分钟电量为统一变量，将储能状态递推、功率限制、容量边界和各类购电费用置于同一计算框架。第一问以LP、MILP和DP交叉核验费用与运行约束；第二至四问按信息实际发布时间生成决策，并以逐区间轨迹复算供需平衡和费用，每项运行约束和费用均可由保存结果核对。

针对不同信息结构，模型采用分层控制。全天计划负责确定普通购电量，储能控制方法处理预测偏差；问题三、问题四-3进一步采用“全日规划、六小时更新”的滚动机制，远期预报用于后续储能规划，每次调整费对应本次正式确认的六小时窗口。问题四-2通过统一年末储电量和连续块重采样比较贪心控制与边际价值DP，问题三、问题四-3采用滚动重建的递减价值DP。

\subsection{模型局限与改进方向}

当前预测只使用历史数据。加入天气、节假日和生产计划等外生变量可改善负载、光伏和实时电价预测；数据分布发生变化时，需要重新校准安全余量与储能策略。问题三采用减购退款并支付50\%违约费的结算解释，交易规则不退回原购电费时，应更换费用函数并重新计算。多日滚动视野和概率预测可用于评估极端天气及价格尖峰下的运行风险。
